% Computer Science A --- Front-end web (HTML, CSS, JavaScript, DOM, async) aligned to \texttt{webcontext.md} lab progression.

\runningheader{Computer Science A}{Web development (HTML, CSS, JS)}{Page \thepage\ of \numpages}

\begingradingrange{csaweb}

\uplevel{\par\textbf{Part 1: HTML \& CSS basics}\par}

\question[4]
What does HTML (HyperText Markup Language) primarily define for a web page?
\begin{choices}
  \choice The visual theme and colors only
  \choice Only interactive behavior and event handlers
  \CorrectChoice The structure and semantic content of the document
  \choice The execution order of JavaScript modules
\end{choices}
\begin{solution}
  \textbf{HTML} marks up headings, paragraphs, lists, links, and other content so the browser can build the DOM tree. \textbf{CSS} handles presentation; \textbf{JavaScript} adds behavior.
\begin{lstlisting}[style=htmlexam]
<!DOCTYPE html>
<html>
  <body>
    <h1>Title</h1>
    <p>Paragraphs and lists carry meaning, not just looks.</p>
  </body>
</html>
\end{lstlisting}
\end{solution}

\question[4]
Which element sets the text typically shown in the browser tab / window title?
\begin{choices}
  \choice \texttt{<h1>}
  \choice \texttt{<header>}
  \CorrectChoice \texttt{<title>}
  \choice \texttt{<meta>}
\end{choices}
\begin{solution}
  The \texttt{<title>} element lives in \texttt{<head>} and names the document. Headings \texttt{<h1>}--\texttt{<h6>} structure visible page content.
\begin{lstlisting}[style=htmlexam]
<head>
  <meta charset="utf-8">
  <title>My lab page</title>
</head>
\end{lstlisting}
\end{solution}

\question[4]
Which fragment correctly represents an \textbf{unordered} list with two items?
\begin{choices}
  \choice \texttt{<ol><li>A</li><li>B</li></ol>}
  \CorrectChoice \texttt{<ul><li>A</li><li>B</li></ul>}
  \choice \texttt{<list><item>A</item></list>}
  \choice \texttt{<ul>A</ul><ul>B</ul>}
\end{choices}
\begin{solution}
  Use \texttt{<ul>} with nested \texttt{<li>} items for bullet lists; \texttt{<ol>} is for numbered lists. Each list item must be an \texttt{<li>}.
\begin{lstlisting}[style=htmlexam]
<ul>
  <li>A</li>
  <li>B</li>
</ul>
\end{lstlisting}
\end{solution}


\newpage

\question[4]
Which element wraps items that should appear as a \textbf{numbered} (ordered) list?
\begin{choices}
  \choice \texttt{<li>}
  \choice \texttt{<dl>}
  \CorrectChoice \texttt{<ol>}
  \choice \texttt{<nav>}
\end{choices}
\begin{solution}
  \texttt{<ol>} (ordered list) works with \texttt{<li>} children for step-like or ranked content; \texttt{<ul>} is for bullets. Definition lists use \texttt{<dl>}/\texttt{<dt>}/\texttt{<dd>}.
\begin{lstlisting}[style=htmlexam]
<ol>
  <li>Step one</li>
  <li>Step two</li>
</ol>
\end{lstlisting}
\end{solution}

\question[4]
In \texttt{<a href="page.html">Open</a>}, what is the role of \texttt{href}?
\begin{choices}
  \CorrectChoice It specifies the link destination (URL or path)
  \choice It sets the font size of the link text
  \choice It makes the link open only in a new window
  \choice It prevents the default navigation behavior
\end{choices}
\begin{solution}
  The \textbf{hypertext reference} attribute \texttt{href} points to the resource to load. Presentation is CSS; \texttt{target} and JS handlers affect navigation behavior separately.
\begin{lstlisting}[style=htmlexam]
<a href="page.html">Open</a>
<a href="https://example.com/" target="_blank">External</a>
\end{lstlisting}
\end{solution}

\question[4]
In the lab progression, \textbf{inline CSS} is applied by placing declarations where?
\begin{choices}
  \choice Only in an external \texttt{.css} file linked via \texttt{<link>}
  \choice Only inside a \texttt{<style>} block in the head
  \CorrectChoice In the element's \texttt{style} attribute as a declaration list
  \choice In a JavaScript \texttt{className} assignment only
\end{choices}
\begin{solution}
  Inline styles use \texttt{style="property: value; ..."} on the element. External and embedded stylesheets are the other main approaches; they cascade with specificity rules.
\begin{lstlisting}[style=htmlexam]
<p style="color: navy; font-size: 14px;">Styled in place</p>
\end{lstlisting}
\end{solution}


\newpage

\question[4]
Which CSS declaration horizontally centers text inside its container (for typical block contexts in the lab)?
\begin{choices}
  \choice \texttt{font-style: italic;}
  \choice \texttt{font-size: 14px;}
  \CorrectChoice \texttt{text-align: center;}
  \choice \texttt{margin: 0 auto;}
\end{choices}
\begin{solution}
  \texttt{text-align} affects inline content within a block. Centering a block itself often uses margins or flexbox; the Homepage/CSS labs emphasized \texttt{text-align} for typography alignment.
\begin{lstlisting}[style=cssexam]
h1 {
  text-align: center;
}
\end{lstlisting}
\end{solution}

\question[4]
The CSS Inline lab emphasizes \texttt{font-style}, \texttt{font-size}, and \texttt{text-align}. Which value of \texttt{font-style} produces italicized text?
\begin{choices}
  \choice \texttt{bold}
  \CorrectChoice \texttt{italic}
  \choice \texttt{underline}
  \choice \texttt{sans-serif}
\end{choices}
\begin{solution}
  \texttt{font-style: italic} slants glyphs; \texttt{font-size} sets length (e.g.\ \texttt{14px}). \texttt{bold} belongs to \texttt{font-weight}; \texttt{underline} is usually \texttt{text-decoration}.
\begin{lstlisting}[style=cssexam]
.note {
  font-style: italic;
  font-size: 14px;
}
\end{lstlisting}
\end{solution}

\question[4]
The rule \texttt{h1 \{ color: blue; \}} applies to which elements?
\begin{choices}
  \CorrectChoice Every \texttt{h1} element (type / element selector)
  \choice Only \texttt{h1} elements with \texttt{id="color"}
  \choice Only \texttt{h1} elements whose \texttt{class} is \texttt{blue}
  \choice The first line of every paragraph
\end{choices}
\begin{solution}
  A bare tag name is an \textbf{element} (type) selector. Combined with classes or IDs it gains specificity; the CSS Selectors lab contrasted element, \texttt{.class}, and \texttt{\#id}.
\begin{lstlisting}[style=cssexam]
h1 {
  color: blue;
}
\end{lstlisting}
\end{solution}

\newpage


\question[4]
In CSS, the selector \texttt{.highlight} targets which elements?
\begin{choices}
  \choice Every element with tag name \texttt{highlight}
  \CorrectChoice Every element whose \texttt{class} attribute includes \texttt{highlight}
  \choice The unique element with \texttt{id="highlight"}
  \choice Only \texttt{<span>} elements
\end{choices}
\begin{solution}
  The \textbf{class selector} prefix \texttt{.} matches the space-separated \texttt{class} list. \textbf{ID} uses \texttt{\#}; \textbf{type} selectors use the bare tag name (\texttt{h1}, etc.).
\begin{lstlisting}[style=cssexam]
.highlight {
  background-color: yellow;
}
\end{lstlisting}
\begin{lstlisting}[style=htmlexam]
<p class="highlight">Important</p>
\end{lstlisting}
\end{solution}

\question[4]
In the CSS box model, which area is \textbf{outside} the border (separating this element from neighbors)?
\begin{choices}
  \choice Padding
  \CorrectChoice Margin
  \choice Content box only
  \choice \texttt{outline} (always identical to border)
\end{choices}
\begin{solution}
  From inside out: \textbf{content} $\rightarrow$ \textbf{padding} $\rightarrow$ \textbf{border} $\rightarrow$ \textbf{margin}. Padding is inside the border; margin is outside.
\begin{lstlisting}[style=cssexam]
.box {
  padding: 8px;
  border: 1px solid #333;
  margin: 16px; /* outside the border */
}
\end{lstlisting}
\end{solution}

\question[4]
In the Box Model lab, \texttt{background-color} applied to an element typically fills which regions (in the standard \texttt{box-sizing: content-box} mental model)?
\begin{choices}
  \choice Only the margin
  \choice Only the content box, excluding padding
  \CorrectChoice Content and padding (under the border edge)
  \choice Only areas outside the border
\end{choices}
\begin{solution}
  Background painting defaults to \textbf{content + padding} up to the inner edge of the border (details vary slightly by shorthand and clipping). The lab tied \texttt{background-color} to the boxed area students style alongside \texttt{padding}, \texttt{margin}, and \texttt{border}.
\begin{lstlisting}[style=cssexam]
.panel {
  padding: 12px;
  background-color: #e8f4ff;
}
\end{lstlisting}
\end{solution}


\newpage
\question[4]
When two rules conflict and both apply to the same element, a rule whose selector is only an \textbf{ID} (\texttt{\#main-title}) will usually beat a rule whose selector is only a \textbf{class} (\texttt{.subtitle}) because:
\begin{choices}
  \choice IDs load later in the file
  \CorrectChoice ID selectors have higher specificity weight than a single class selector
  \choice Classes are deprecated in modern CSS
  \choice The browser always prefers the shorter rule
\end{choices}
\begin{solution}
  Specificity approximates \textbf{inline} $>$ \textbf{ID} $>$ \textbf{class/attribute} $>$ \textbf{type}. Tie-breaking uses source order and \texttt{!important} (sparingly). The lab table summarized ID vs class vs element priority.
\begin{lstlisting}[style=cssexam]
.subtitle { color: gray; }
#main-title { color: navy; /* wins over .subtitle */ }
\end{lstlisting}
\end{solution}

\question[4]
The line \texttt{<meta charset="utf-8">} in \texttt{<head>} primarily tells the browser:
\begin{choices}
  \choice Which JavaScript file to load first
  \choice The exact width of the layout viewport
  \CorrectChoice How to interpret the document's bytes as characters (text encoding)
  \choice That the page must not be cached
\end{choices}
\begin{solution}
  The charset declaration names an encoding (here UTF-8) so the parser maps downloaded bytes to the right Unicode characters---especially important for accents and symbols.
\begin{lstlisting}[style=htmlexam]
<head>
  <meta charset="utf-8">
  <title>Encoding demo</title>
</head>
\end{lstlisting}
\end{solution}

\question[4]
Which markup correctly links an \textbf{external} stylesheet \texttt{styles.css} from the same folder?
\begin{choices}
  \choice \texttt{<script href="styles.css"></script>}
  \CorrectChoice \texttt{<link rel="stylesheet" href="styles.css">}
  \choice \texttt{<style src="styles.css"></style>}
  \choice \texttt{<css href="styles.css" />}
\end{choices}
\begin{solution}
  \texttt{<link rel="stylesheet" href="...">} pulls a \texttt{.css} file into the cascade. \texttt{<style>} holds inline CSS text; \texttt{<script>} is for scripts.
\begin{lstlisting}[style=htmlexam]
<head>
  <link rel="stylesheet" href="styles.css">
</head>
\end{lstlisting}
\end{solution}

\newpage

\question[4]
In HTML, \texttt{class="card featured"} on one element means:
\begin{choices}
  \choice The element has invalid markup (only one class name is allowed)
  \choice \texttt{featured} is a tag name nested inside \texttt{card}
  \CorrectChoice The element has \textbf{both} classes \texttt{card} and \texttt{featured} (space-separated list)
  \choice \texttt{card} is the ID and \texttt{featured} is the tag
\end{choices}
\begin{solution}
  The \texttt{class} attribute is a space-separated token list. Selectors \texttt{.card}, \texttt{.featured}, and \texttt{.card.featured} can all match the same node.
\begin{lstlisting}[style=htmlexam]
<article class="card featured">Sale item</article>
\end{lstlisting}
\begin{lstlisting}[style=cssexam]
.card { border: 1px solid #ccc; }
.featured { background: #fff8e6; }
\end{lstlisting}
\end{solution}

\question[4]
In CSS, the selector \texttt{*} (asterisk alone) means:
\begin{choices}
  \choice Only elements named \texttt{star}
  \choice Only the first child of every parent
  \CorrectChoice All elements in the document (universal selector)
  \choice A syntax error in modern CSS
\end{choices}
\begin{solution}
  The universal selector matches every element; it is often used in resets or broad rules (use carefully because it can hurt performance and specificity intuition).
\begin{lstlisting}[style=cssexam]
* { box-sizing: border-box; }
\end{lstlisting}
\end{solution}

\question[4]
Without extra rules, which property is typically \textbf{inherited} from a parent element to its descendants?
\begin{choices}
  \choice \texttt{margin-top}
  \choice \texttt{border-width}
  \CorrectChoice \texttt{color}
  \choice \texttt{width}
\end{choices}
\begin{solution}
  Many text-oriented properties (\texttt{color}, \texttt{font-family}, \texttt{font-size} in the normal case) inherit so children match the parent's typography. Box properties like \texttt{margin} and \texttt{border} usually do not inherit.
\begin{lstlisting}[style=htmlexam]
<section style="color: navy;">
  <p>Inherited text color.</p>
</section>
\end{lstlisting}
\end{solution}

\newpage

\question[4]
The \texttt{alt} attribute on \texttt{<img>} is important mainly because it:
\begin{choices}
  \choice Sets the image file path (same role as \texttt{src})
  \choice Forces the image to load faster
  \CorrectChoice Provides a text substitute for the image (accessibility, broken image, screen readers)
  \choice Defines the clickable region for the image only
\end{choices}
\begin{solution}
  \texttt{alt} describes the image content for assistive tech and when the asset fails to load; \texttt{src} is the actual URL/path.
\begin{lstlisting}[style=htmlexam]
<img src="logo.png" alt="Company logo">
\end{lstlisting}
\end{solution}

\question[4]
In HTML5, \texttt{<br>} and \texttt{<img>} are usually written without a closing tag because they are:
\begin{choices}
  \choice Deprecated and should be replaced by \texttt{<div>}
  \choice Only valid inside \texttt{<table>}
  \CorrectChoice Void (empty) elements---they do not wrap text content
  \choice Required to use self-closing XML form \texttt{</img>} in all browsers
\end{choices}
\begin{solution}
  Void elements have no element content; the slash-only XHTML-style \texttt{<br />} is tolerated but optional in HTML5 for these tags.
\begin{lstlisting}[style=htmlexam]
<p>Line one<br>Line two</p>
<img src="pic.png" alt="">
\end{lstlisting}
\end{solution}

\question[4]
Which statement best describes \textbf{semantic} HTML elements?
\begin{choices}
  \choice A generic \texttt{<div>} is the main example of a semantic element
  \choice Semantic tags are unrelated to document structure
  \CorrectChoice They name meaningful regions (e.g.\ \texttt{<header>}, \texttt{<nav>}, \texttt{<article>}) so structure and grouping carry meaning
  \choice They are developer-only comments that browsers ignore
\end{choices}
\begin{solution}
  Semantic elements describe roles (navigation, main content, aside) rather than being purely presentational wrappers. \texttt{<div>} is non-semantic until given a role or ARIA.
\end{solution}

\question[4]
What is the HTML5 doctype declaration?
\begin{choices}
  \CorrectChoice \texttt{<!DOCTYPE html>}
  \choice \texttt{<!DOCTYPE HTML PUBLIC "-//W3C//DTD HTML 4.01//EN" "http://www.w3.org/TR/html4/strict.dtd">}
  \choice \texttt{<!HTML doctype>}
  \choice \texttt{<DOCTYPE html>}
\end{choices}
\begin{solution}
  HTML5 uses the short, case-insensitive \texttt{<!DOCTYPE html>} so browsers use standards mode.
\begin{lstlisting}[style=htmlexam]
<!DOCTYPE html>
<html lang="en">
  <head><meta charset="utf-8"><title>Demo</title></head>
  <body></body>
</html>
\end{lstlisting}
\end{solution}

\question[4]
Which of these is \textbf{not} a valid CSS unit?
\begin{choices}
  \choice \texttt{\%} (percentage)
  \choice \texttt{em} (font-relative)
  \CorrectChoice \texttt{li} (not a unit---it names a list item in HTML)
  \choice \texttt{px} or \texttt{pt}
\end{choices}
\begin{solution}
  Common units include \texttt{px}, \texttt{em}, \texttt{rem}, \texttt{\%}, \texttt{vh}/\texttt{vw}, \texttt{pt}, etc. \texttt{li} is an element, not a length unit.
\end{solution}

\question[4]
HTML is a programming language.
\begin{choices}
  \choice True
  \CorrectChoice False---HTML is a \textbf{markup} language for structure and semantics, not a general-purpose programming language
  \choice True only when scripts are disabled
  \choice True only for static pages
\end{choices}
\begin{solution}
  Behavior and computation belong to languages like JavaScript; HTML declares content and structure.
\end{solution}

\question[4]
Which group lists elements that normally belong inside \texttt{<body>}?
\begin{choices}
  \CorrectChoice \texttt{<div>}, \texttt{<header>}, \texttt{<span>}, \texttt{<strong>}
  \choice \texttt{<link>}, \texttt{<meta>}, \texttt{<p>}, \texttt{<title>}
  \choice \texttt{<head>}, \texttt{<h1>}, \texttt{<h2>}, \texttt{<h3>} (only in \texttt{<head>})
  \choice \texttt{<article>}, \texttt{<aside>}, \texttt{<nav>}, \texttt{<title>}
\end{choices}
\begin{solution}
  \texttt{<title>}, \texttt{<meta>}, and \texttt{<link>} (for metadata/stylesheets) live in \texttt{<head>}. Visible flow content usually sits in \texttt{<body>}.
\end{solution}

\question[4]
Which values are \textbf{valid} for an \texttt{<input>} element's \texttt{type} attribute?
\begin{choices}
  \CorrectChoice \texttt{password} and \texttt{text} are both valid \texttt{type} values
  \choice \texttt{title} and \texttt{name} are valid \texttt{type} values
  \choice Only \texttt{name} is valid; \texttt{text} is not
  \choice Only \texttt{password} exists in HTML5
\end{choices}
\begin{solution}
  \texttt{name} and \texttt{title} are attributes, not \texttt{type} keywords. Many input types exist (\texttt{text}, \texttt{password}, \texttt{email}, \texttt{checkbox}, etc.).
\begin{lstlisting}[style=htmlexam]
<input type="text" name="user">
<input type="password" name="pass">
\end{lstlisting}
\end{solution}

\question[4]
Which of the following is \textbf{not} an HTML5 \emph{element} name?
\begin{choices}
  \choice \texttt{<body>}
  \choice \texttt{<h1>}
  \choice \texttt{<p>}
  \CorrectChoice \texttt{<class>} (\texttt{class} is an \emph{attribute}, not a tag)
\end{choices}
\begin{solution}
  Use \texttt{class="..."} on elements; there is no \texttt{<class>} tag.
\end{solution}

\question[4]
Order the CSS box model layers from \textbf{outermost} to \textbf{innermost}.
\begin{choices}
  \choice border $\rightarrow$ margin $\rightarrow$ content $\rightarrow$ padding
  \choice margin $\rightarrow$ padding $\rightarrow$ border $\rightarrow$ content
  \CorrectChoice margin $\rightarrow$ border $\rightarrow$ padding $\rightarrow$ content
  \choice border $\rightarrow$ content $\rightarrow$ margin $\rightarrow$ padding
\end{choices}
\begin{solution}
  From outside in: \textbf{margin}, \textbf{border}, \textbf{padding}, \textbf{content}. This matches the mental model used alongside \texttt{background-color} painting under the border.
\end{solution}

\question[4]
Which element holds an \textbf{embedded} (internal) stylesheet in the document?
\begin{choices}
  \choice \texttt{<script>}
  \choice \texttt{<body>}
  \CorrectChoice \texttt{<style>} (usually in \texttt{<head>})
  \choice \texttt{<css>}
\end{choices}
\begin{solution}
  Put CSS rules inside \texttt{<style>}; link external files with \texttt{<link rel="stylesheet" href="...">}.
\begin{lstlisting}[style=htmlexam]
<head>
  <style>
    body { font-family: system-ui, sans-serif; }
  </style>
</head>
\end{lstlisting}
\end{solution}

\question[4]
What is true of a typical \textbf{inline} element (e.g.\ \texttt{<span>}) in normal flow?
\begin{choices}
  \CorrectChoice It sits in the text line by default; horizontal size follows content (subject to wrapping)
  \choice It always starts on a new line like a block box
  \choice It always stretches to the full viewport width
  \choice The browser truncates its text when width is tight
\end{choices}
\begin{solution}
  Inline boxes participate in line layout; block-level boxes (e.g.\ \texttt{<div>}) typically break onto new lines and can take full width when set to \texttt{display: block}.
\end{solution}

\newpage
\uplevel{\par\textbf{Part 2: JavaScript fundamentals}\par}

\question[4]
You need a counter that increments on each button click. Which declaration is most appropriate for the counter binding if it will be reassigned?
\begin{choices}
  \choice \texttt{const count = 0;} (reassign later anyway)
  \CorrectChoice \texttt{let count = 0;}
  \choice \texttt{var count = 0;} only---\texttt{let} cannot be used for counters
  \choice \texttt{function count = 0;}
\end{choices}
\begin{solution}
  \texttt{let} is block-scoped and allows reassignment. \texttt{const} cannot be rebound; the bound value may still be mutable if it is an object. \texttt{var} is legacy-scoped.
\begin{lstlisting}[style=jsexam]
let count = 0;
btn.onclick = () => { count += 1; };
\end{lstlisting}
\end{solution}

\question[4]
In the If-statement lab style, an integer \texttt{n} is \textbf{even} when which expression is true?
\begin{choices}
  \choice \texttt{n / 2 === 0}
  \CorrectChoice \texttt{n \% 2 === 0}
  \choice \texttt{n \% 2 === 1}
  \choice \texttt{Math.floor(n) === n / 2}
\end{choices}
\begin{solution}
  The \textbf{remainder} operator \texttt{\%} yields zero for even \texttt{n} under \texttt{n \% 2 === 0}. Odd numbers leave remainder \texttt{1}.
\begin{lstlisting}[style=jsexam]
if (n % 2 === 0) {
  console.log("even");
}
\end{lstlisting}
\end{solution}

\question[4]
In \texttt{if (cond1) \{ A \} else if (cond2) \{ B \} else \{ C \}}, block \texttt{C} runs when:
\begin{choices}
  \choice \texttt{cond1} and \texttt{cond2} are both truthy
  \CorrectChoice \texttt{cond1} is falsy and \texttt{cond2} is falsy
  \choice \texttt{cond1} is truthy regardless of \texttt{cond2}
  \choice Only when \texttt{cond2} throws an error
\end{choices}
\begin{solution}
  The chain tests \texttt{cond1} then, if it failed, \texttt{cond2}. The \texttt{else} branch runs only when every tested condition above was falsy (If-statement lab structure).
\begin{lstlisting}[style=jsexam]
if (cond1) {
  A();
} else if (cond2) {
  B();
} else {
  C();
}
\end{lstlisting}
\end{solution}

\newpage

\question[4]
In \texttt{for (init; condition; final)}, the \texttt{final} clause (for example \texttt{i++}) runs:
\begin{choices}
  \choice Once before any iteration
  \choice After the body on each iteration
  \choice Only if the loop body throws
  \CorrectChoice After each iteration body completes, before the next condition check
\end{choices}
\begin{solution}
  A C-style \texttt{for} loop runs \textbf{init} once, then each cycle: test \textbf{condition}, run body, run \textbf{final}, repeat. That matches the lab's three-part loop table.
\begin{lstlisting}[style=jsexam]
for (let i = 0; i < 3; i++) {
  console.log(i);
}
\end{lstlisting}
\end{solution}

\question[4]
Which value is \textbf{falsy} in JavaScript when used as an \texttt{if} condition?
\begin{choices}
  \choice \texttt{[]} (empty array)
  \choice \texttt{\{\}} (empty object)
  \CorrectChoice \texttt{""} (empty string)
  \choice \texttt{"0"}
\end{choices}
\begin{solution}
  The short falsy list includes \texttt{false}, \texttt{0}, \texttt{-0}, \texttt{0n}, \texttt{""}, \texttt{null}, \texttt{undefined}, \texttt{NaN}. Empty arrays and objects are \textbf{truthy}---a common beginner surprise from the Truthyfalsy lab.
\begin{lstlisting}[style=jsexam]
if ("") { /* skipped */ }
if ([]) { /* runs: [] is truthy */ }
\end{lstlisting}
\end{solution}

\question[4]
What is the result of \texttt{5 === "5"} in JavaScript?
\begin{choices}
  \CorrectChoice \texttt{false} (no type coercion with strict equality)
  \choice \texttt{true} because both represent five
  \choice \texttt{undefined}
  \choice A syntax error
\end{choices}
\begin{solution}
  \textbf{Strict equality} \texttt{===} compares \textbf{without} coercion, so number and string differ. \texttt{5 == "5"} is \texttt{true} with loose equality (coercion). Modern style prefers \texttt{===} / \texttt{!==}.
\begin{lstlisting}[style=jsexam]
console.log(5 === "5"); // false
\end{lstlisting}
\end{solution}


\newpage
\question[4]
What is the value of \texttt{5 == "5"} (loose equality) in JavaScript?
\begin{choices}
  \CorrectChoice \texttt{true} (operands may be coerced before comparison)
  \choice \texttt{false}
  \choice \texttt{"55"}
  \choice It throws \texttt{TypeError}
\end{choices}
\begin{solution}
  Loose \texttt{==} applies the language's coercion rules; numeric string \texttt{"5"} converts to number \texttt{5}. The Typecoercion lab contrasts this with predictable \texttt{===}.
\begin{lstlisting}[style=jsexam]
console.log(5 == "5"); // true
\end{lstlisting}
\end{solution}

\question[4]
What does \texttt{Number("42")} evaluate to?
\begin{choices}
  \CorrectChoice The number \texttt{42}
  \choice The string \texttt{"42"}
  \choice \texttt{NaN}
  \choice \texttt{true}
\end{choices}
\begin{solution}
  \texttt{Number(...)} performs \textbf{explicit} numeric conversion. Invalid strings yield \texttt{NaN}; the lab contrasted this with implicit coercion in \texttt{==} comparisons.
\begin{lstlisting}[style=jsexam]
const x = Number("42"); // 42
\end{lstlisting}
\end{solution}

\question[4]
For an array \texttt{arr}, \texttt{arr.push("x")} primarily:
\begin{choices}
  \choice Removes the last element
  \CorrectChoice Adds an element at the end and returns the new \texttt{length}
  \choice Sorts the array in place
  \choice Returns the pushed value only, without changing \texttt{arr}
\end{choices}
\begin{solution}
  \texttt{push} appends; \texttt{pop} removes from the end. \texttt{length} tracks element count (for dense arrays, last index is \texttt{length - 1}).
\begin{lstlisting}[style=jsexam]
const arr = [1, 2];
arr.push("x"); // length is now 3
\end{lstlisting}
\end{solution}

\question[4]
Given \texttt{let user = \{ name: "Ada", score: 10 \};}, which expression reads the score using \textbf{bracket} notation with a string key?
\begin{choices}
  \choice \texttt{user.score.}
  \CorrectChoice \texttt{user["score"]}
  \choice \texttt{user::score}
  \choice \texttt{user(score)}
\end{choices}
\begin{solution}
  \textbf{Dot} syntax (\texttt{user.score}) is idiomatic for valid identifiers. \textbf{Brackets} support dynamic keys and non-identifier keys (\texttt{user["score"]}).
\begin{lstlisting}[style=jsexam]
const user = { name: "Ada", score: 10 };
const s = user["score"];
\end{lstlisting}
\end{solution}

\question[4]
In the Button-click lab style, which API retrieves a single element by its \texttt{id} attribute?
\begin{choices}
  \choice \texttt{document.querySelector("id")}
  \CorrectChoice \texttt{document.getElementById("...")}
  \choice \texttt{document.getElementsByTagName("id")}
  \choice \texttt{window.findId("...")}
\end{choices}
\begin{solution}
  \texttt{getElementById} returns one \texttt{Element} or \texttt{null}. The lab pairs this with \texttt{onclick} assignment or later \texttt{addEventListener} patterns.
\begin{lstlisting}[style=jsexam]
const btn = document.getElementById("go");
\end{lstlisting}
\end{solution}

\question[4]
To replace all text inside a DOM element \texttt{el} with the string \texttt{"Done"}, you should assign to:
\begin{choices}
  \choice \texttt{el.innerHTML = "Done"} only (always safest)
  \CorrectChoice \texttt{el.textContent = "Done"} (plain text, avoids parsing HTML)
  \choice \texttt{el.value = "Done"} on any element
  \choice \texttt{el.node = "Done"}
\end{choices}
\begin{solution}
  \texttt{textContent} sets character data without interpreting markup. \texttt{innerHTML} parses HTML (XSS risk with untrusted strings). \texttt{.value} applies to form controls.
\begin{lstlisting}[style=jsexam]
el.textContent = "Done";
\end{lstlisting}
\end{solution}

\question[4]
To read what the user currently typed in a text \texttt{<input>} element referenced as \texttt{inputEl}, use:
\begin{choices}
  \choice \texttt{inputEl.textContent}
  \CorrectChoice \texttt{inputEl.value}
  \choice \texttt{inputEl.getAttribute("value")} always matches live typing
  \choice \texttt{inputEl.innerText} only
\end{choices}
\begin{solution}
  The live string for most inputs is \texttt{.value}. The HTML \texttt{value} attribute is only the \emph{initial} default unless you sync it; the Input lab stresses reading \texttt{.value} for current text.
\begin{lstlisting}[style=jsexam]
const text = inputEl.value;
\end{lstlisting}
\end{solution}

\question[4]
What is the value of \texttt{typeof 42} in JavaScript?
\begin{choices}
  \choice \texttt{"integer"}
  \choice \texttt{"float"}
  \CorrectChoice \texttt{"number"}
  \choice \texttt{"decimal"}
\end{choices}
\begin{solution}
  JavaScript has one numeric type in \texttt{typeof} results: \texttt{"number"} (both \texttt{3} and \texttt{3.14}).
\begin{lstlisting}[style=jsexam]
typeof 42;        // "number"
typeof "42";       // "string"
\end{lstlisting}
\end{solution}

\question[4]
What does \texttt{"5" + 1} evaluate to?
\begin{choices}
  \choice \texttt{6} (numeric addition)
  \choice \texttt{NaN}
  \CorrectChoice \texttt{"51"} (string concatenation after coercion)
  \choice The string \texttt{"6"}
\end{choices}
\begin{solution}
  With \texttt{+}, if either operand is a string, the other is typically coerced to string and concatenated---not added numerically.
\begin{lstlisting}[style=jsexam]
console.log("5" + 1);   // "51"
console.log("5" - 1);   // 4  (subtraction uses numeric coercion)
\end{lstlisting}
\end{solution}

\question[4]
Which expression uses a \textbf{template literal} to embed \texttt{name} into a greeting?
\begin{choices}
  \choice \texttt{'Hello, ' + name + '!'}
  \CorrectChoice \texttt{`Hello, \$\{name\}!`}
  \choice \texttt{"Hello, \#name!"}
  \choice \texttt{concat("Hello, ", name)}
\end{choices}
\begin{solution}
  Backtick strings allow \texttt{\$\{\ldots\}} interpolations; they are the ES6 template literal syntax.
\begin{lstlisting}[style=jsexam]
const name = "Ada";
const msg = `Hello, ${name}!`;
\end{lstlisting}
\end{solution}

\question[4]
\texttt{JSON.parse('\{"score":10\}')} returns:
\begin{choices}
  \choice A string equal to \texttt{"score"}
  \CorrectChoice A JavaScript object (e.g.\ you can read \texttt{.score})
  \choice \texttt{undefined}
  \choice A syntax error because single quotes are illegal
\end{choices}
\begin{solution}
  \texttt{JSON.parse} deserializes a JSON text into JS values (objects, arrays, strings, numbers, booleans, null). \texttt{JSON.stringify} does the reverse.
\begin{lstlisting}[style=jsexam]
const obj = JSON.parse('{"score":10,"ok":true}');
console.log(obj.score);
\end{lstlisting}
\end{solution}

\question[4]
Compared to \texttt{splice}, \texttt{Array.prototype.slice} is different because \texttt{slice}:
\begin{choices}
  \choice Always empties the original array
  \choice Only works on strings, not arrays
  \CorrectChoice Returns a shallow copy of a section and does \textbf{not} mutate the original array
  \choice Sorts the array in descending order
\end{choices}
\begin{solution}
  \texttt{slice(start, end)} copies a range; \texttt{splice} changes the array in place (insert/remove). Recognizing mutating vs non-mutating methods prevents subtle bugs.
\begin{lstlisting}[style=jsexam]
const nums = [10, 20, 30];
const part = nums.slice(0, 2); // [10, 20], nums unchanged
\end{lstlisting}
\end{solution}

\question[4]
The operator \texttt{!==} is best described as:
\begin{choices}
  \choice Loose inequality (allows type coercion)
  \CorrectChoice Strict inequality (no type coercion, opposite of \texttt{===})
  \choice Bitwise NOT combined with assignment
  \choice Legal only inside \texttt{for} loops
\end{choices}
\begin{solution}
  Pair \texttt{!==} with \texttt{===} for predictable comparisons; \texttt{!=} is the loose form paired with \texttt{==}.
\begin{lstlisting}[style=jsexam]
5 !== "5"; // true
5 !=  "5"; // false
\end{lstlisting}
\end{solution}

\question[4]
For a non-empty array \texttt{arr}, the index of the \textbf{last} element is:
\begin{choices}
  \choice \texttt{arr.length}
  \CorrectChoice \texttt{arr.length - 1}
  \choice \texttt{arr.length + 1}
  \choice Always \texttt{0}
\end{choices}
\begin{solution}
  Arrays are zero-indexed: first at \texttt{0}, last at \texttt{length - 1}. Off-by-one errors at \texttt{length} are a common debugging target.
\begin{lstlisting}[style=jsexam]
const arr = ["a", "b", "c"];
const last = arr[arr.length - 1];
\end{lstlisting}
\end{solution}

\question[4]
What are the results of \texttt{3 == '3'} and \texttt{3 === '3'} in JavaScript?
\begin{choices}
  \CorrectChoice \texttt{true} then \texttt{false}---\texttt{==} may coerce types; \texttt{===} does not
  \choice \texttt{false} then \texttt{true}
  \choice Both are \texttt{false}
  \choice Both are \texttt{true}
\end{choices}
\begin{solution}
  Loose equality coerces the string \texttt{"3"} to the number \texttt{3}. Strict equality compares type and value, so number and string differ.
\begin{lstlisting}[style=jsexam]
console.log(3 == "3");  // true
console.log(3 === "3"); // false
\end{lstlisting}
\end{solution}

\question[4]
Which keywords can declare a variable binding in modern JavaScript?
\begin{choices}
  \CorrectChoice \texttt{let}, \texttt{var}, and \texttt{const}
  \choice \texttt{int} and \texttt{async}
  \choice Only \texttt{var} (\texttt{let} and \texttt{const} are invalid)
  \choice \texttt{for} and \texttt{function} only
\end{choices}
\begin{solution}
  \texttt{let} and \texttt{const} are block-scoped; \texttt{var} is function-scoped (legacy). \texttt{const} forbids rebinding but object contents may still mutate.
\end{solution}

\question[4]
JavaScript (ES6) has separate built-in \texttt{int} and \texttt{double} numeric types.
\begin{choices}
  \choice True
  \CorrectChoice False---there is one \texttt{number} type (IEEE-754); \texttt{typeof} reports \texttt{"number"}
  \choice True only in strict mode
  \choice True only in browsers, not in Node
\end{choices}
\begin{solution}
  All numeric literals share the same type; distinguish integers with \texttt{Number.isInteger} when needed.
\end{solution}

\question[4]
Which line correctly creates an empty \textbf{object literal}?
\begin{choices}
  \CorrectChoice \texttt{let myObject = \{\};}
  \choice \texttt{let Object myObject = new Object();}
  \choice \texttt{let myObject = ();}
  \choice JavaScript has no object type
\end{choices}
\begin{solution}
  Curly braces with optional properties define object literals; \texttt{new Object()} works but the literal form is idiomatic.
\begin{lstlisting}[style=jsexam]
let myObject = {};
myObject.score = 10;
\end{lstlisting}
\end{solution}

\question[4]
How do you read the \textbf{first} element of an array \texttt{arr}?
\begin{choices}
  \choice \texttt{arr(0)}
  \CorrectChoice \texttt{arr[0]}
  \choice \texttt{arr[1]}
  \choice \texttt{arr\{0\}}
\end{choices}
\begin{solution}
  Arrays are zero-indexed: first at \texttt{0}, last at \texttt{length - 1}.
\end{solution}

\question[4]
You need a template literal that interpolates \texttt{product.quantity}, \texttt{product.name}, and \texttt{product.price}, and prints a literal dollar sign immediately before the numeric price. Which pattern is correct?
\begin{choices}
  \choice Backticks with \texttt{\$product.quantity} (no braces)
  \CorrectChoice Backticks with \texttt{\$\{product.quantity\}}, \texttt{\$\{product.name\}}, and \texttt{\$\$\{product.price\}} (first \texttt{\$} is literal, \texttt{\$\{...\}} interpolates)
  \choice Single quotes with \texttt{\$\{\{product.quantity\}\}}
  \choice Single quotes with \texttt{\$\{product.quantity\}}
\end{choices}
\begin{solution}
  Inside template literals, \texttt{\$\{expr\}} evaluates an expression. To output a literal \texttt{\$}, escape by writing \texttt{\$\$} before \texttt{\$\{product.price\}}.
\begin{lstlisting}[style=jsexam]
const s = `${product.quantity} units of ${product.name} at $${product.price}`;
\end{lstlisting}
\end{solution}

\question[4]
What happens when this code runs?
\begin{lstlisting}[style=jsexam]
const a = 0;
for (let i = 5; i < 10; i++) {
  a++;
}
\end{lstlisting}
\begin{choices}
  \choice \texttt{a} is \texttt{5}
  \choice \texttt{a} is \texttt{4}
  \choice \texttt{a} is \texttt{undefined}
  \CorrectChoice A \textbf{TypeError}: \texttt{const} bindings cannot be reassigned (\texttt{a++} mutates the binding)
\end{choices}
\begin{solution}
  \texttt{const} fixes the binding to the value \texttt{0}; increment operators require a mutable binding---use \texttt{let} for counters.
\end{solution}

\newpage
\uplevel{\par\textbf{Part 3: DOM manipulation \& events}\par}

\question[4]
To change which image file an \texttt{<img>} displays, you can assign to \texttt{imgEl.src} or call:
\begin{choices}
  \choice \texttt{imgEl.href = "pic.png"}
  \CorrectChoice \texttt{imgEl.setAttribute("src", "pic.png")}
  \choice \texttt{imgEl.alt = "pic.png"}
  \choice \texttt{imgEl.createElement("src")}
\end{choices}
\begin{solution}
  The Attributes lab uses \texttt{.src} or \texttt{setAttribute("src", ...)} so the browser loads a new resource. \texttt{alt} is alternative text, not the image URL.
\begin{lstlisting}[style=jsexam]
imgEl.setAttribute("src", "pic.png");
/* or: imgEl.src = "pic.png"; */
\end{lstlisting}
\end{solution}

\question[4]
To add a new \texttt{<li>} to an existing \texttt{<ol>} in the Appending-elements lab pattern, you typically:
\begin{choices}
  \choice Call \texttt{ol.innerHTML += "<li>"} only (no creation step)
  \CorrectChoice \texttt{document.createElement("li")} then \texttt{parent.appendChild(newLi)}
  \choice \texttt{ol.push(li)}
  \choice \texttt{new Element("li")}
\end{choices}
\begin{solution}
  \texttt{createElement} builds a node; \texttt{appendChild} attaches it under a parent. String \texttt{innerHTML} reparses markup and can drop event listeners---\texttt{createElement} is the structured approach.
\begin{lstlisting}[style=jsexam]
const li = document.createElement("li");
li.textContent = "New item";
ol.appendChild(li);
\end{lstlisting}
\end{solution}

\question[4]
Compared to assigning \texttt{element.onclick = fn}, \texttt{element.addEventListener("click", fn)} is especially useful because it:
\begin{choices}
  \choice Runs the handler before the event exists
  \CorrectChoice Can register \textbf{multiple} listeners for the same event without overwriting earlier ones
  \choice Only works for keyboard events
  \choice Prevents event bubbling always
\end{choices}
\begin{solution}
  \texttt{onclick} property assignment replaces the previous handler. \texttt{addEventListener} stacks listeners and offers capture/bubbling control (\texttt{Eventlistener} lab).
\begin{lstlisting}[style=jsexam]
el.addEventListener("click", fn1);
el.addEventListener("click", fn2);
\end{lstlisting}
\end{solution}

\newpage

\question[4]
Inside a click handler, \texttt{event.target} typically refers to:
\begin{choices}
  \choice The \texttt{window} object always
  \CorrectChoice The innermost element that was the origin of the event
  \choice The document root always
  \choice The function that handled the event
\end{choices}
\begin{solution}
  The \textbf{Event} object reports \texttt{type}, \texttt{target} (where the event started), and related fields. The Eventobject lab used \texttt{target} after a click.
\begin{lstlisting}[style=jsexam]
el.addEventListener("click", (event) => {
  console.log(event.target);
});
\end{lstlisting}
\end{solution}

\question[4]
\texttt{document.querySelector("\#go")} differs from \texttt{document.getElementById("go")} mainly in that \texttt{querySelector}:
\begin{choices}
  \choice Can only match tags, never IDs
  \CorrectChoice Takes a full \textbf{CSS selector string} and returns the \textbf{first} matching element (or \texttt{null})
  \choice Always returns a NodeList of all matches
  \choice Throws if the ID does not exist
\end{choices}
\begin{solution}
  \texttt{querySelector} is the general CSS-selector API; prefix IDs with \texttt{\#}, classes with \texttt{.}, etc. \texttt{getElementById} is ID-only and very slightly older/simpler.
\begin{lstlisting}[style=jsexam]
const a = document.querySelector("#go");
const b = document.getElementById("go");
\end{lstlisting}
\end{solution}

\question[4]
Calling \texttt{event.preventDefault()} inside a click handler on a link typically:
\begin{choices}
  \choice Deletes the link from the DOM
  \CorrectChoice Stops the browser's default action (e.g.\ following the \texttt{href}) while still allowing the handler to run
  \choice Removes all other listeners on that element
  \choice Is required for every \texttt{addEventListener} call
\end{choices}
\begin{solution}
  \texttt{preventDefault} blocks default behavior (navigation, form submit, context menu, etc.). It is unrelated to \texttt{stopPropagation}, which affects event bubbling/capture.
\begin{lstlisting}[style=jsexam]
link.addEventListener("click", (event) => {
  event.preventDefault();
  console.log("No navigation");
});
\end{lstlisting}
\end{solution}

\newpage
\question[4]
To add a CSS class \texttt{active} without overwriting existing classes on \texttt{el}, the preferred DOM API is:
\begin{choices}
  \choice \texttt{el.class = "active"}
  \choice \texttt{el.className += "active"} only (no space)
  \CorrectChoice \texttt{el.classList.add("active")}
  \choice \texttt{el.setAttribute("classList", "active")}
\end{choices}
\begin{solution}
  \texttt{classList} maintains the class token set: \texttt{add}, \texttt{remove}, \texttt{toggle}, \texttt{contains}. Assigning \texttt{className} replaces the whole attribute string.
\begin{lstlisting}[style=jsexam]
el.classList.add("active");
el.classList.remove("hidden");
\end{lstlisting}
\end{solution}

\question[4]
For \texttt{<div id="app" data-user-id="7">}, the JavaScript property that reads \texttt{"7"} is:
\begin{choices}
  \choice \texttt{el.data-user-id} (valid property syntax)
  \choice \texttt{el.getAttribute("user-id")}
  \CorrectChoice \texttt{el.dataset.userId}
  \choice \texttt{el.dataset.data-user-id}
\end{choices}
\begin{solution}
  \texttt{data-*} attributes appear on \texttt{element.dataset} in camelCase: \texttt{data-user-id} $\rightarrow$ \texttt{dataset.userId}. Values are always strings.
\begin{lstlisting}[style=htmlexam]
<div id="app" data-user-id="7" data-role="admin"></div>
\end{lstlisting}
\begin{lstlisting}[style=jsexam]
const el = document.getElementById("app");
console.log(el.dataset.userId); // "7"
\end{lstlisting}
\end{solution}

\question[4]
To remove a child node \texttt{child} from its parent \texttt{parent}, modern code often uses:
\begin{choices}
  \choice \texttt{parent.delete(child)}
  \choice \texttt{child.erase()}
  \CorrectChoice \texttt{child.remove()} (or \texttt{parent.removeChild(child)})
  \choice \texttt{document.unlink(child)}
\end{choices}
\begin{solution}
  \texttt{Element.remove()} detaches the node from the tree. The older \texttt{removeChild} pattern still appears in legacy code but is more verbose.
\begin{lstlisting}[style=jsexam]
child.remove();
/* or: parent.removeChild(child); */
\end{lstlisting}
\end{solution}

\newpage

\uplevel{\par\textbf{Part 4: Advanced JavaScript}\par}

\question[4]
In an ES6 class, which special method runs when you write \texttt{new MyClass(args)}?
\begin{choices}
  \choice \texttt{static init()}
  \CorrectChoice \texttt{constructor(args) \{ \ldots \}}
  \choice \texttt{function MyClass()}
  \choice \texttt{this.build()}
\end{choices}
\begin{solution}
  \texttt{constructor} initializes the instance; \texttt{this} refers to that instance inside instance methods. The Classes lab covered instantiation and calling methods like \texttt{getDescription()}.
\begin{lstlisting}[style=jsexam]
class Widget {
  constructor(name) { this.name = name; }
}
const w = new Widget("clock");
\end{lstlisting}
\end{solution}

\question[4]
Which construct runs \texttt{recover()} only if \texttt{risky()} throws?
\begin{choices}
  \choice \texttt{if (risky()) recover();}
  \CorrectChoice \texttt{try \{ risky(); \} catch (e) \{ recover(); \}}
  \choice \texttt{throw risky() || recover();}
  \choice \texttt{risky().catch(recover)} for sync code only
\end{choices}
\begin{solution}
  \texttt{try}/\texttt{catch} handles synchronous exceptions; \texttt{throw} signals an error. The Errors lab stressed preventing uncaught exceptions from crashing the script path under test.
\begin{lstlisting}[style=jsexam]
try { risky(); } catch (e) { recover(); }
\end{lstlisting}
\end{solution}

\question[4]
\texttt{setTimeout(callback, 1000)} schedules the \texttt{callback} to run:
\begin{choices}
  \choice Immediately in the same call stack
  \CorrectChoice After roughly 1000\,ms, via the browser's timer queue (async relative to current synchronous code)
  \choice Only if \texttt{callback} returns \texttt{true}
  \choice Before any \texttt{console.log} on the page
\end{choices}
\begin{solution}
  \textbf{Callbacks} passed to \texttt{setTimeout} run later; order relative to other macrotasks depends on the event loop (Callbackfunction lab).
\begin{lstlisting}[style=jsexam]
setTimeout(() => console.log("later"), 1000);
\end{lstlisting}
\end{solution}

\newpage

\question[4]
Declaring a function as \texttt{async function f() \{ \ldots \}} means that calling \texttt{f()} returns:
\begin{choices}
  \choice A plain number or string always
  \choice \texttt{undefined} always
  \CorrectChoice A \textbf{Promise} (resolved with the function's returned value or rejected on throw)
  \choice A synchronous blocking lock
\end{choices}
\begin{solution}
  \texttt{async} functions wrap the return value in a Promise. \texttt{await} pauses \emph{inside} the async function until a Promise settles (Asyncawait lab).
\begin{lstlisting}[style=jsexam]
async function f() { return 1; }
const p = f(); // a Promise
\end{lstlisting}
\end{solution}

\question[4]
The expression \texttt{await fetch(url)} (inside an \texttt{async} function) typically yields:
\begin{choices}
  \choice Parsed JSON immediately
  \CorrectChoice A \texttt{Response} object (body still unread until \texttt{json()}, \texttt{text()}, etc.)
  \choice A string URL
  \choice An \texttt{ArrayBuffer} always
\end{choices}
\begin{solution}
  \texttt{fetch} resolves to a \texttt{Response}; you then \texttt{await res.json()} or similar to parse. The Fetch lab chained request $\rightarrow$ parse $\rightarrow$ DOM update.
\begin{lstlisting}[style=jsexam]
async function load(url) {
  const res = await fetch(url);
  const data = await res.json();
  return data;
}
\end{lstlisting}
\end{solution}

\question[4]
In \texttt{Math.max(...values)}, the \texttt{...values} syntax in the call is called:
\begin{choices}
  \choice A rest parameter
  \CorrectChoice Spread (expands an iterable into individual arguments)
  \choice A template literal
  \choice Destructuring assignment
\end{choices}
\begin{solution}
  \textbf{Spread} in a call expands an array (or iterable) into arguments. \textbf{Rest} in the parameter list (\texttt{function f(...args)}) collects extra arguments into an array (Spread/rest lab).
\begin{lstlisting}[style=jsexam]
const values = [3, 9, 4];
Math.max(...values);
\end{lstlisting}
\end{solution}


\newpage
\question[4]
In \texttt{function sum(...nums) \{ return nums.reduce((a,b)=>a+b,0); \}}, the \texttt{...nums} in the parameter list is:
\begin{choices}
  \choice Spread syntax
  \CorrectChoice Rest syntax (collects remaining arguments into an array \texttt{nums})
  \choice An object spread
  \choice Illegal in strict mode
\end{choices}
\begin{solution}
  Rest parameters bundle trailing arguments. Spread and rest share the \texttt{...} token but appear in opposite ``directions'' (call vs.\ parameter list).
\begin{lstlisting}[style=jsexam]
function sum(...nums) {
  return nums.reduce((a, b) => a + b, 0);
}
\end{lstlisting}
\end{solution}

\question[4]
After \texttt{const res = await fetch(url);}, a robust pattern is to check \texttt{res.ok} because:
\begin{choices}
  \choice \texttt{fetch} throws automatically on HTTP 404
  \CorrectChoice \texttt{fetch} resolves even for error status codes (4xx/5xx); \texttt{res.ok} is false when the status is not in the 200--299 range
  \choice \texttt{res.ok} is true only before \texttt{await}
  \choice \texttt{res.ok} replays the \texttt{network\_request}
\end{choices}
\begin{solution}
  Failed HTTP responses still yield a \texttt{Response}; inspect \texttt{status} or the boolean \texttt{ok} before calling \texttt{json()}.
\begin{lstlisting}[style=jsexam]
const res = await fetch(url);
if (!res.ok) throw new Error(`HTTP ${res.status}`);
const data = await res.json();
\end{lstlisting}
\end{solution}

\question[4]
If an \texttt{async} function body executes \texttt{throw new Error("bad");}, the Promise returned by calling that function will:
\begin{choices}
  \choice Stay pending forever
  \choice Always resolve to \texttt{undefined}
  \CorrectChoice Become \textbf{rejected} with that error (unless caught inside the function)
  \choice Throw synchronously in the caller without a Promise
\end{choices}
\begin{solution}
  Exceptions inside \texttt{async} functions reject the implicit Promise. Callers should use \texttt{try}/\texttt{catch} with \texttt{await} or \texttt{.catch} on the Promise.
\begin{lstlisting}[style=jsexam]
async function boom() { throw new Error("bad"); }
boom().catch((e) => console.log(e.message));
\end{lstlisting}
\end{solution}

\newpage

\question[4]
\texttt{Promise.all([p1, p2])} resolves when:
\begin{choices}
  \choice The \textbf{first} promise settles (either fulfilled or rejected)
  \CorrectChoice \textbf{Every} promise in the array has fulfilled; if any rejects, the whole \texttt{Promise.all} rejects
  \choice All promises are still pending
  \choice Only when called inside a \texttt{for} loop
\end{choices}
\begin{solution}
  \texttt{Promise.all} is useful for parallel async work; combine with \texttt{Promise.allSettled} when you need all outcomes regardless of failures (ES2020+).
\begin{lstlisting}[style=jsexam]
const [a, b] = await Promise.all([fetch(url1), fetch(url2)]);
\end{lstlisting}
\end{solution}

\question[4]
The method \texttt{res.json()} on a \texttt{fetch} \texttt{Response} is \textbf{asynchronous} because:
\begin{choices}
  \choice It always returns \texttt{undefined}
  \choice It requires two URLs
  \CorrectChoice Reading/parsing the body is streamed and completes later---it returns a \textbf{Promise}
  \choice JSON is only valid in Node.js, not browsers
\end{choices}
\begin{solution}
  Body consumption methods (\texttt{json}, \texttt{text}, \texttt{blob}) are async; you must \texttt{await} (or use \texttt{.then}) even after you already have \texttt{res}.
\begin{lstlisting}[style=jsexam]
const res = await fetch(url);
const data = await res.json();
\end{lstlisting}
\end{solution}

\question[4]
Given:
\begin{lstlisting}[style=jsexam]
let promise = new Promise((resolve, reject) => {
  reject("Promise rejected!");
});
promise.catch((error) => console.log(error));
promise.catch((error) => console.log(error));
\end{lstlisting}
What is printed?
\begin{choices}
  \choice Nothing
  \CorrectChoice \texttt{Promise rejected!} appears \textbf{twice} (each \texttt{.catch} on the same promise can handle the rejection)
  \choice \texttt{undefined} twice
  \choice Only an unhandled rejection warning
\end{choices}
\begin{solution}
  Registering multiple \texttt{.catch} handlers on one promise means each handler may run for the same rejection (they are separate registrations on the same settlement).
\end{solution}

\question[4]
In JavaScript, function values are objects and inherit (via the prototype chain) from \texttt{Object}.
\begin{choices}
  \CorrectChoice True
  \choice False
  \choice True only for arrow functions
  \choice True only for class methods
\end{choices}
\begin{solution}
  Functions are callable objects; their prototypes ultimately link like other objects, which is why you can attach properties to a function.
\end{solution}

\question[4]
What is \textbf{event capturing} in the DOM event model?
\begin{choices}
  \CorrectChoice The event is dispatched from the root down through ancestors toward the target (\textbf{capture} phase) before the target and bubble phases
  \choice The event propagates only from the target upward (bubbling only)
  \choice The browser logs every event to the console automatically
  \choice Events jump between sibling nodes at random
\end{choices}
\begin{solution}
  With \texttt{addEventListener(..., true)}, handlers run during capture; the default bubble phase goes target $\rightarrow$ ancestors. \texttt{stopPropagation} affects both.
\end{solution}

\question[4]
Which kind of JavaScript value cannot be represented in standard JSON text?
\begin{choices}
  \CorrectChoice A function value
  \choice A string
  \choice A plain object
  \choice An array
\end{choices}
\begin{solution}
  JSON supports strings, numbers, booleans, \texttt{null}, arrays, and objects. Functions, \texttt{undefined}, and symbols are omitted or require custom serializers.
\end{solution}

\question[4]
Which API serializes a JavaScript value into a JSON \textbf{string}?
\begin{choices}
  \choice \texttt{JSON.toString()}
  \CorrectChoice \texttt{JSON.stringify()}
  \choice \texttt{String.JSON()}
  \choice \texttt{json.toString()}
\end{choices}
\begin{solution}
  Pair \texttt{JSON.stringify} with \texttt{JSON.parse} for round-trips; handle replacers/revivers when dates or special values matter.
\begin{lstlisting}[style=jsexam]
const s = JSON.stringify({ score: 10 });
const obj = JSON.parse(s);
\end{lstlisting}
\end{solution}

\question[4]
The browser \texttt{fetch} API is primarily used to \ldots
\begin{choices}
  \choice Read object fields without referencing the object
  \choice Force any Promise to resolve synchronously
  \CorrectChoice Perform HTTP requests and obtain \texttt{Response} objects (often followed by \texttt{json()} or \texttt{text()})
  \choice Search lexical scope for a variable name
\end{choices}
\begin{solution}
  \texttt{fetch} is the modern network entry point; it returns a Promise of a \texttt{Response} and does not throw on HTTP error status codes by default.
\end{solution}

\endgradingrange{csaweb}
